\section{探索问题的本质}

\subsection{为什么探索很重要}

在强化学习中，智能体必须在\textbf{利用}（exploitation，选择已知最好的动作）和\textbf{探索}（exploration，尝试可能更好的动作）之间取得平衡。

\begin{figure}[H]
\centering
\includegraphics[width=0.9\textwidth]{slides-images/slide-04.jpg}
\caption{探索难题的直觉展示：有些任务对人类简单，对从零开始的 RL 却几乎不可做}
\end{figure}

\begin{importantbox}{探索的核心困境}
如果我们总是选择当前估值最高的动作（贪心策略），可能永远发现不了真正的最优策略。例如在多臂赌博机中，如果第一次拉臂恰好获得了正奖励，贪心策略会一直拉同一个臂，即使其他臂的期望奖励更高。
\end{importantbox}

\subsection{为什么 Montezuma's Revenge 成了经典探索基准}

课程先用 Montezuma's Revenge 这种稀疏奖励游戏说明问题：拿到钥匙有奖励、开门有奖励、被 skull 杀死却未必立刻给出足够明确的负反馈，而完成任务的关键行为之间间隔又非常长。人类之所以知道要做什么，不是因为看到了 reward，而是因为理解了 sprite 的语义和任务结构。

\begin{figure}[H]
\centering
\includegraphics[width=0.9\textwidth]{slides-images/slide-05.jpg}
\caption{Montezuma's Revenge：奖励稀疏、行为长链、语义理解缺失导致探索困难}
\end{figure}

\subsection{把自己放到算法的位置上}

讲者进一步用牌类游戏 Mao 做类比：你只知道 ``违反规则会受罚''，但规则本身只能靠试错发现，而且规则未必直观。这恰好刻画了长时序稀疏奖励任务的本质困难。

\begin{figure}[H]
\centering
\includegraphics[width=0.9\textwidth]{slides-images/slide-06.jpg}
\caption{Mao 类比：规则只能通过试错发现，任务越长、规则越隐蔽，探索越困难}
\end{figure}

\subsection{探索和利用其实是同一个问题}

课程强调，两种常见表述本质等价：
\begin{itemize}
    \item ``如何发现需要复杂长链行为才能到达的高奖励策略？''
    \item ``何时该继续做当前已知最好的事，何时该尝试新行为？''
\end{itemize}
前者强调 temporally extended behavior，后者强调 exploration-exploitation trade-off，但它们都在问同一件事：\textbf{今天愿不愿意承受一点损失，去换取明天可能更高的回报。}

\begin{figure}[H]
\centering
\includegraphics[width=0.9\textwidth]{slides-images/slide-07.jpg}
\caption{餐馆、广告投放、油井勘探：探索与利用的日常例子}
\end{figure}

\begin{warningbox}{探索之所以难，不只是因为动作空间大}
真正棘手的地方是：很多关键探索行为在局部上看不到收益，甚至会让短期回报下降。于是如果训练信号只盯着即时表现，agent 就很容易被推向保守但次优的局部策略。
\end{warningbox}

\begin{figure}[H]
\centering
\includegraphics[width=0.9\textwidth]{slides-images/slide-08.jpg}
\caption{从 bandit 到大规模无结构 MDP：探索问题的理论难度迅速上升}
\end{figure}

\subsection{探索在 RL 应用中的体现}

\begin{itemize}
    \item \textbf{机器人}：探索不同的抓取策略、运动模式
    \item \textbf{大语言模型}：在 RLHF/RL for reasoning 中，探索不同的推理路径
    \item \textbf{推荐系统}：探索用户可能喜欢但尚未展示过的内容
\end{itemize}

\subsection{本章小结}

探索是所有 RL 问题的根本挑战之一。没有有效探索，智能体会陷入次优策略。

